%%%%%%%%%%%%%%%%%%%%%%%%%%%%%%%%%%%%%%%%%%%%%%%%%%%%%%%%%%%%%%%%%%%%%%
% Lớp 10 | Véctơ (chưa xét tọa độ) | Tích vô hướng (chưa xét tọa độ)
% Dạng 0H5?4-6: Cực trị và chứng minh bất đẳng thức
% Số câu: 1  (sinh tự động bởi tach_cau.py; từ nay thầy sửa trực tiếp file này)
%%%%%%%%%%%%%%%%%%%%%%%%%%%%%%%%%%%%%%%%%%%%%%%%%%%%%%%%%%%%%%%%%%%%%%

%%%%%%%%%%%%%%%% TỰ LUẬN %%%%%%%%%%%%%%%%

% nguồn: MaHoaNH_TuLuan_Lop10
\begin{bt}%[0H5V4-6]%[ma:d4f3f54c]
	Cho đoạn $AB=20$. Tồn tại điểm $M$ sao cho $T=3MA^2+2 MB^2$ đạt giá trị bé nhất  $T_{\min}$. Tính giá trị $T_{\min }$.
	\loigiai{ 
		Gọi điểm  $I$ thỏa mãn $3 \overrightarrow{IA}+2 \overrightarrow{IB}=\overrightarrow 0$. Khi đó\\
		\[3\overrightarrow{IA}+2\left(\overrightarrow{IA}+\overrightarrow{AB}\right)=\overrightarrow 0\Leftrightarrow 5\overrightarrow{IA}+2\overrightarrow{AB}=\overrightarrow  0\Leftrightarrow \overrightarrow{AI}=\dfrac{2}{5}\overrightarrow{AB}.\]
		Suy ra điểm $I$ thuộc đoạn $AB$  và $IA=\dfrac{2}{5} \cdot AB=\dfrac{2}{5} \cdot 20=8, IB=12$.\\
		Ta  có
		\begin{eqnarray*}
			T&=&3 MA^2+2  MB^2\\
			&=&3 \overrightarrow{MA}^2+2 \overrightarrow{MB}^2\\ 
			 &=&3(\overrightarrow{MI}+\overrightarrow{IA})^2+2\left(\overrightarrow{MI}+\overrightarrow{IB}\right)^2\\
			&=&3{{\overrightarrow{MI}}^2}+6\overrightarrow{MI}\cdot \overrightarrow{IA}+3{{\overrightarrow{IA}}^2}+2{{\overrightarrow{MI}}^2}+4\overrightarrow{MI}\cdot  \overrightarrow{IB}+2{{\overrightarrow{IB}}^2} \\
			&=&5MI^2+3Ia^2+2IB^2+2\overrightarrow{MI}\left(\underbrace{3\overrightarrow{IA}+2\overrightarrow{IB}}_{\overrightarrow  0}\right)\\
			 &=&5MI^2+3IA^2+2IB^2.
		\end{eqnarray*}
		Ta có $\left(3IA^2+2IB^2\right)$ là hằng số do ba điểm $A$,  $B$, $I$ cố định.\\
		Do đó $T$ đạt giá trị nhỏ nhất $\Leftrightarrow 5MI^2$ nhỏ nhất $\Leftrightarrow MI$ bé nhất $\Leftrightarrow$ Điểm $M$ trùng với điểm  $I$.\\
		Khi đó giá trị $T$ nhỏ nhất là $T_{\min}=3IA^2+2 IB^2=3 \cdot 8^2+2 \cdot 12^2=480$. 
	} 
\end{bt}
